\section{Ruling out star height one}\label{section: lower bound}

We prove \Cref{theorem: lower bound}.
Let $E$ be a term with $\sh(E) \le 1$.
Every star in $E$ has the form $F^{*}$ with $F$ star-free.
Let $q$ be at least $\cd(F)$ for all subterms $F^{*}$ of $E$.
Take the modules $\M^{\pm} = \M^{\pm}_q$ of \Cref{proposition: local pair} and the cell models $\C(w)$ of \Cref{definition: cells}.

\subparagraph*{Coordinates.}
Let $\mathbb{V}$ be the common vertex set of $\M^{+}$ and $\M^{-}$, and let $\Theta \defeq \mathbb{V} \uplus \set{\pi_0, \pi_1}$ be the set of \emph{tags}.
For $w = w_1 \cdots w_n$, let $\hat w \defeq {\triangleright}\, w$ be the word over $\Lambda \defeq \set{{\triangleright}, 0, 1}$ with positions $0, \dots, n$.
We call $\set{p^{0}_0, p^{0}_1}$ \emph{block $0$}.
For $i \ge 1$, \emph{block $i$} is $\range{u_i, v_i} \cup \set{p^{i}_0, p^{i}_1}$.
A vertex of $\C(w)$ in block $i$ gets the coordinates $(i, \vartheta)$.
Here $\vartheta \in \mathbb{V}$ is its vertex in the module, or $\pi_j$ if it is the port $p^{i}_j$.
Order $\Theta$ by the order of $\mathbb{V}$ followed by $\pi_0 < \pi_1$.
Then the order of $\C(w)$ is the lexicographic order of coordinates.

\begin{lemma}[Normalization]\label{lemma: normalization}
  Let $F^{*}$ be a subterm of $E$, and let $x = (i, \vartheta)$ and $y = (j, \vartheta')$ be vertices of $\C(w)$ with $x \le y$.
  Then the truth value of $F^{*}(x, y)$ in $\C(w)$ depends only on $\vartheta, \vartheta', \hat w_i, \hat w_j$ and $j - i$.
\end{lemma}
\begin{proof}
  Let $w'$ arise from $w$ by setting $w_k \defeq 0$ for every $k$ with $\range{u_k, v_k} \subseteq \range{x, y}$.
  We pass from $\C(w)$ to $\C(w')$ one module at a time.
  Each step replaces a copy of $\M^{+}$ by a copy of $\M^{-}$ in an interval $\range{u_k, v_k} \subseteq \range{x, y}$.
  The conditions of \Cref{lemma: replacement} hold.
  (\ref{item: same outside}) holds since the modules differ only in a $\vc$-pair with an interior target.
  (\ref{item: isolated}) holds since no input pair leaves a module.
  (\ref{item: local equivalence}) holds by \Cref{proposition: local pair}.
  Moreover, $x$ and $y$ are outer.
  So $F^{*}(x, y)$ has the same value in $\C(w)$ and in $\C(w')$.
  In $\C(w')$ all modules strictly between blocks $i$ and $j$ are copies of $\M^{-}$.
  So by \Cref{lemma: locality}, the value is determined by $\vartheta, \vartheta'$, $w'_i, w'_j$ (when $i, j \ge 1$) and $j - i$.
  These are in turn determined by $\vartheta, \vartheta', \hat w_i, \hat w_j$ and $j - i$.
\end{proof}

For a subterm $F^{*}$ of $E$, let $f_F(\vartheta, a, \vartheta', a', k) \in \set{0,1}$ be the value given by \Cref{lemma: normalization}, for $a, a' \in \Lambda$ and $k \ge 0$.

\begin{lemma}[Periodicity]\label{lemma: periodicity}
  There are $N_0 \in \nat$ and $d \ge 1$ such that $f_F(\vartheta, a, \vartheta', a', k) = f_F(\vartheta, a, \vartheta', a', k + d)$ for all subterms $F^{*}$ of $E$, all $\vartheta, a, \vartheta', a'$ and all $k \ge N_0$.
\end{lemma}
\begin{proof}
  Since there are finitely many choices of $F, \vartheta, a, \vartheta', a'$, it suffices to show that each sequence $k \mapsto f_F(\vartheta, a, \vartheta', a', k)$ is ultimately periodic.
  For $k \ge 1$, the normalized interval $\range{x, y}$ in the proof of \Cref{lemma: normalization} is an ordered sum $X \oplus B^{k-1} \oplus Y$.
  Here $X$ is the rest of block $i$ from $x$, $B$ is a copy of $\M^{-}$ followed by two ports, and $Y$ is the part of block $j$ up to $y$.
  No input pair connects different summands.
  Write the vertices of $X$, $B$ and $Y$ as letters of a finite alphabet, using disjoint copies.
  Then $X \oplus B^{k-1} \oplus Y$ is first-order interpretable in the word that lists its vertices.
  Two positions lie in the same summand iff their distance equals the distance of their letters inside that summand, and this is first-order expressible with a bounded number of successor steps.
  Since $F^{*}$ is definable in monadic second-order logic, the set of $k$ for which $F^{*}(\min, \max)$ holds is the preimage of a regular language under $k \mapsto \delta_X \delta_B^{k-1} \delta_Y$~\cite{buchiWeakSecondorderArithmetic1960}.
  Hence it is ultimately periodic.
\end{proof}

\subparagraph*{Translation into first-order logic with modular predicates.}
Let $\FO[<, \MOD{d}]$ be first-order logic over words in $\Lambda^{*}$.
It has the order, a unary predicate $P_a$ for each $a \in \Lambda$, and unary predicates $\mathrm{mod}^{r}_{d}(i)$, meaning $i \equiv r \pmod d$, for $r \in \range{0, d-1}$.
Expressions such as $j = i + k$ for a constant $k$, $j \ge i + N_0$, and $j - i \equiv r \pmod d$ are definable.
The last one is $\bigvee_{r'} \mathrm{mod}^{r'}_d(i) \wedge \mathrm{mod}^{r' + r}_d(j)$.

\begin{lemma}[Translation]\label{lemma: translation}
  For every subterm $G$ of $E$ and all tags $\vartheta, \vartheta' \in \Theta$, there is an $\FO[<, \MOD{d}]$ formula $\varphi^{\vartheta, \vartheta'}_{G}(i, j)$ with the following property.
  For all $w \in \set{0,1}^{+}$ and all vertices $(i, \vartheta) \le (j, \vartheta')$ of $\C(w)$, we have $G((i, \vartheta), (j, \vartheta'))$ in $\C(w)$ iff $\hat w \models \varphi^{\vartheta, \vartheta'}_{G}(i, j)$.
\end{lemma}
\begin{proof}
  By induction on $G$.
  Let $\mathrm{vert}_{\vartheta}(i)$ express that $(i, \vartheta)$ is a vertex, that is, $i > 0$ or $\vartheta \in \set{\pi_0, \pi_1}$.
  Let $\mathrm{le}_{\vartheta,\vartheta'}(i, j)$ express $(i, \vartheta) \le (j, \vartheta')$.
  \begin{itemize}
    \item Variables $o \in \set{\vr, \vc, \vb}$: every input pair lies in one module, so we take $i = j \wedge \bigvee \set{P_a(i) \mid a \in \set{0,1},\ (\vartheta, \vartheta') \in o^{\M^{(a)}}}$.
          The constants $\id$ and $\emp$ are similar.
    \item $\com{G}$: take $\mathrm{le}_{\vartheta,\vartheta'}(i, j) \wedge \neg\varphi^{\vartheta, \vartheta'}_{G}$, and $G_1 \union G_2$: take the disjunction.
    \item $G_1 \compo G_2$: take $\exists k \bigvee_{\vartheta''} (\mathrm{vert}_{\vartheta''}(k) \wedge \varphi^{\vartheta, \vartheta''}_{G_1}(i, k) \wedge \varphi^{\vartheta'', \vartheta'}_{G_2}(k, j))$.
    \item $F^{*}$: by \Cref{lemma: normalization,lemma: periodicity}, take the disjunction over $a, a' \in \Lambda$ of $P_a(i) \wedge P_{a'}(j) \wedge \psi_{a, a'}(i, j)$.
          Here $\psi_{a, a'}$ is the disjunction of the following formulas.
          First, $j = i + k$ for each $k < N_0$ with $f_F(\vartheta, a, \vartheta', a', k) = 1$.
          Second, $j \ge i + N_0 \wedge j - i \equiv r \pmod d$ for each $r$ such that $f_F(\vartheta, a, \vartheta', a', k) = 1$ for the $k \ge N_0$ with $k \equiv r$.
          \qedhere
  \end{itemize}
\end{proof}

\begin{lemma}[Parity is not definable]\label{lemma: parity undefinable}
  No $\FO[<, \MOD{d}]$ sentence $\varphi$ satisfies, for all $w \in \set{0,1}^{+}$: ${\triangleright} w \models \varphi$ iff $w$ has an odd number of $1$s.
\end{lemma}
\begin{proof}
  Restrict to $w = (1\, 0^{d-1})^{n}$, which has $n$ ones.
  The word ${\triangleright} w$ has positions $0, \dots, nd$.
  It is first-order interpretable in the linear order $\range{0, n}$ with $d$ copies.
  Position $0$ is $(0, 0)$, and position $(k-1)d + r + 1$ is $(k, r)$ for $k \in \range{n}$ and $r \in \range{0, d-1}$.
  The order is lexicographic.
  $P_{\triangleright}$ holds iff $k = 0$, $P_1$ iff $k \ne 0$ and $r = 0$, and $P_0$ iff $k \ne 0$ and $r \ne 0$.
  $\mathrm{mod}^{r'}_d$ holds at $(k, r)$ with $k \ne 0$ iff $r + 1 \equiv r' \pmod d$.
  All of these are first-order definable in $(\range{0, n}, <)$.
  Hence $\varphi$ would yield a first-order sentence $\varphi^{\flat}$ such that $(\range{0, n}, <) \models \varphi^{\flat}$ iff $n$ is odd.
  But linear orders of size at least $2^{\ell}$ are indistinguishable by sentences of quantifier rank $\ell$~\cite{ebbinghausFiniteModelTheory1995}.
\end{proof}

\begin{proof}[Proof of \Cref{theorem: lower bound}]
  Assume that $E$ is equivalent to $\Psit$ on $\Cls{2}$.
  Every $\C(w)$ lies in $\Cls{2}$ by \Cref{lemma: cells pathwidth}, and its relations are irreflexive.
  By \Cref{lemma: parity}, $E((0, \pi_0), (n, \pi_1))$ holds in $\C(w)$ iff $w$ has an odd number of $1$s.
  By \Cref{lemma: translation}, this property is defined by the sentence $\exists i\, \exists j\, (i \text{ is first} \wedge j \text{ is last} \wedge \varphi^{\pi_0, \pi_1}_{E}(i, j))$, contradicting \Cref{lemma: parity undefinable}.
\end{proof}

\begin{remark}\label{remark: more variables}
  The proof does not use that $E$ is a term over $\vr, \vc, \vb$ only.
  If $E$ contains further variables, we interpret them as empty in $\C(w)$.
  This changes neither the graph nor any argument above.
  So $\Psit$ is not equivalent on $\Cls{2}$ to any term of star height at most one, over any finite set of variables.
\end{remark}
